Another way to solve a system with unsymmetric \(A\), is with the Petrov-Galerkin approach, \cite{Vorst03}, which is to find \(x_m\) such that the residual \(r_m\) is orthogonal to another \(m\)-dimensional subspace. The choice for this method is subspace is the Krylov subspace generated by \(A^T\) and \(r_0\), \(\mathcal{K}^m(A^T;r_0)\), which will lead to the iterative solver known as Bi-Conjugate Gradient (Bi-CG). An elegant way to derive the recurrence relations for this method is to use CG on this specific system, \cite{Vorst03},
\begin{equation}
    \begin{pmatrix}
    0 & A \\ A^T & 0
    \end{pmatrix}
    \begin{pmatrix}
    \hat{x} \\ x
    \end{pmatrix}
    = \begin{pmatrix}
    b \\ \hat{b}
    \end{pmatrix},
\end{equation}
with the preconditioner
\begin{equation}
    \hat{M}=\begin{pmatrix}
    0 & M \\ M^T & 0
    \end{pmatrix},
\end{equation}
written as \(\hat{M}^{-1}B\tilde{x}=\hat{M}^{-1}\tilde{b}\).

With this construction CG can be used to solve the system since \(\hat{M}^{-1}B\) is symmetric for the inner product \([x,y]:=(x,My)\), which is not a proper inner product since \(\hat{M}^{-1}B\) is not positive definite, \cite{Vorst03}. Because it is not a inner product this shows a problem with Bi-CG, as this can be evaluated to zero when a prober inner product wouldn't, this is known as breakdown of the method.

\noindent Which gets us the relations,
\begin{align*}
    p_i&=M^{-1}r_{i-1}+\beta_{i-1}p_{u-1}\\
    \tilde{p}_i&=M^{-T}\tilde{r}_{i-1}+\beta_{i-1}\tilde{p}_{u-1}\\
    x_i&=x_{i-1}+\alpha_ip_i\\
    r_i&=r_{i-1}-\alpha_iAp_i\\
    \tilde{r}_i&=\tilde{r}_{i-1}-\alpha_iA^T\tilde{p}_i\\
\intertext{with scalers}
    \alpha_i&=\frac{\rho_{i-1}}{\tilde{p}^T_i A p_i} \\
    \beta_i&=\frac{\rho_{i-1}}{rho_{i-2}}\\
    \rho_{i-1}&=(M^{-T}\tilde{r}_{i-1})^TM^{-1}r_{i-1}
\end{align*}
with initial residual \(r_0=b-Ax_0\) and \(\tilde{r}_0\) such that \((r_0, \tilde{r}_0)\ne0\), often chosen to be \(\tilde{r}_0=r_0\). With the bi-orthogonality and bi-conjugate, between the two sets of vectors, \cite{Saad2003},
\begin{align}
    (r_j, \tilde{r}_i)&=0, \text{ for } j\ne i\label{eq: BiCG orthogonal}\\
    (Ap_j, \tilde{p}_i)&=0, \text{ for } j\ne i,
\end{align}

\noindent ensures that Bi-CG would converge in exact arithmetic. As mentioned breakdown can occur with Bi-CG, other disadvantages with this method is the extra matrix vector products and as with using the system of normal equations is the need to be able to calculated the transpose of the original systems matrix, as said the difficulty depends on the implementation of the matrix. A more serious problem with Bi-CG and the methods derived from it, is that in implementation the methods suffer from irregular convergence. This means that norm of the current residual may be larger than the previous residual, even if the algorithm does converges.